\documentclass[a4paper,14pt]{extarticle}
\usepackage[utf8]{inputenc}
\usepackage[english,russian]{babel}

\usepackage{amsthm}
\usepackage{graphicx}
\usepackage{caption}
\usepackage{amssymb}
\usepackage{amsmath}
\usepackage{mathrsfs}
\usepackage{euscript}
\usepackage{graphicx}
\usepackage{subfig}
\usepackage{caption}
\usepackage{color}
\usepackage{bm}
\usepackage{tabularx}
\usepackage{adjustbox}
\usepackage[toc,page]{appendix}
\usepackage{comment}
\usepackage{rotating}


% Определение новых математических функций 
\DeclareMathOperator{\const}{const}
\DeclareMathOperator{\Lap}{Lap}
\DeclareMathOperator{\softmax}{softmax}
\DeclareMathOperator{\KL}{KL}
\DeclareMathOperator{\argmin}{arg \hspace{3pt} min}


\newtheorem{theorem}{Теорема}
\newtheorem{lemma}[theorem]{Лемма}
\newtheorem{definition}{Определение}[section]

\numberwithin{equation}{section}

\newcommand*{\No}{No.}

\begin{document}

% Титульный лист
%\input{./frontpage.tex}

% Аннотация
%\newpage


\begin{abstract}
Работа посвящена оптимизации структуры нейронной сети. Предполагается, что число параметров нейросети можно существенно снизить без значимой потери качества
и значимого повышения дисперсии функции ошибки. Предлагается метод прореживания параметров нейронной сети, основанный на автоматическом определении релевантности параметров. Для определения релевантности параметров предлагается проанализировать ковариационую матрицу апостериорного распределения параметров и удалить из нейросети
мультикоррелирующие параметры. Для определения мультикорреляции используется метод Белсли. Для анализа качества представленного алгоритма проводятся эксперименты на выборке Boston Housing, а также на синтетических данных.



\smallskip
\textbf{Ключевые слова}: нейронные сети; оптимизация гиперпараметров; метод Белсли; релевантность параметров; прореживание нейронной сети
\end{abstract}






% Нумерация должна начинаться со второй страницы
\setcounter{page}{2}

% Оглавление
\newpage
\tableofcontents

% Обозначения и сокращения
% \input{./dict.tex}

% Введение
%\newpage


\section{Введение}
Решается задача выбора оптимальной структуры нейронной сети. В силу высокой вычислительной сложности, время оптимизации нейронных сетей может занимать до нескольких дней. Поэтому построение и выбороптимальной структуры нейронной сети также является вычислительно сложной процедурой, которая значимо влияет на итоговое качество
модели. Использование избыточно сложных моделей с избыточным числом неинформативных параметров является препятствием для использования глубоких сетей на мобильных устройствах в режиме реального времени. Существует ряд подходов к построению оптимальной сети. В работах [2, 3] предлагается использовать модель градиентного спуска для оптимизации сети. В [4] используются байесовские методы оптимизации параметров нейронных сетей. Другим методом поиска оптимальной структуры является прореживание избыточно сложной модели.
В работе [6] предлагается удалять наименее релевантные параметры на основе значений первой и второй производных функции ошибки.
Данная работа посвящена прореживанию структуры сети. Предлагается удалять наименее релевантные параметры модели. Под релевантностью [6] подразумевается то, насколько параметр влияет на функцию
ошибки. Малая релевантность указывает на то, что удаление этого параметра не влечет значимого изменения функции ошибки. Метод предлагает построение исходной избыточной сложности нейросети с большим числом избыточных параметров. Для определения релевантности параметров предлагается оптимизацировать параметры и гиперпараметры в
единой процедуре. Для удаления параметров предлагается использовать метод Белсли.
Проверка и анализ метода проводится на выборке Boston Housing, Wine и синтетических данных. Результат сравнивается с моделью, полученной при помощи базовых алгоритмов.

\section{Анализ литературы}
To be done...


% Постановка задачи
\newpage


\section{Постановка задачи}
Задана выборка 
\begin{equation} \label{task.1} 
\mathfrak{D} = \{\mathbf{x}_i, y_i\}, \quad i = 1, \ldots, N,
\end{equation}
где $\mathbf{x}_i \in \mathbb{R}^m$, $y_i \in \{1, \ldots, Y\}$, $Y$ -- число классов. Рассмотрим модель $f(\mathbf{x}, \mathbf{w}): \mathbb{R}^m \times \mathbb{R}^n \to \{1, \ldots, Y\}$, где $\mathbf{w} \in \mathbb{R}^n$ -- пространство параметров модели,
\begin{equation} \label{task.2} 
f(\mathbf{x}, \mathbf{w}) = \softmax \big(f_1(f_2(\ldots(f_l(\mathbf{x}, \mathbf{w})\big),
\end{equation}
где $f_k(\mathbf{x}, \mathbf{w}) = \tanh(\mathbf{w}^{\mathsf{T}}\mathbf{x})$, $l$ -- число слоев нейронной сети, $k \in \{1, \ldots, l\}$. 

Параметр $w_i$ модели $f$ называется активным, если $w_i \neq 0$. Множество индексов активных параметров обозначим $\mathcal{A} \subset \mathcal{J} = \{1, \ldots, n\}$. Задано пространство параметров модели:
\begin{equation} \label{task.3} 
\mathbb{W}_{\mathcal{A}} = \{\mathbf{w} \in \mathbb{R}^n \mid w_j \neq 0, \hspace{5pt} j \in \mathcal{A}\}.
\end{equation}

Для модели $f$ с множеством индексов активных параметров $\mathcal{A}$ и соответствующего ей вектора параметров $\mathbf{w} \in \mathbb{W}_{\mathcal{A}}$ определим логарифмическую функцию правдоподобия выборки:
\begin{equation} \label{task.4} 
\mathcal{L}_{\mathfrak{D}}(\mathfrak{D}, \mathcal{A}, \mathbf{w}) = \log p(\mathfrak{D} \mid \mathcal{A}, \mathbf{w}), 
\end{equation}
где $p(\mathfrak{D} \mid \mathcal{A}, \mathbf{w})$ -- апостериорная вероятность выборки $\mathfrak{D}$ при заданных $\mathbf{w}$, $\mathbf{A}$. Оптимальные значения $\mathbf{w}$, $\mathbf{A}$ находятся из минимизации $-\mathcal{L}_{\mathcal{A}}(\mathfrak{D}, \mathcal{A})$ -- логарифма правдоподобия модели:
\begin{equation} \label{task.5} 
\mathcal{L}_{\mathcal{A}}(\mathfrak{D}, \mathcal{A}) = \log p(\mathfrak{D} \mid \mathcal{A}) = \log \int_{\mathbf{w \in \mathbb{W}_{\mathcal{J}}}} p(\mathfrak{D} \mid \mathbf{w})p(\mathbf{w} \mid \mathcal{A}) d \mathbf{w},
\end{equation}
где $p(\mathbf{w} \mid \mathcal{A})$ -- априорная вероятность вектора параметров в пространстве $\mathbb{W}_{\mathcal{J}}$.

Так как вычисление интеграла \ref{task.5} является вычислительно сложной задачей, рассмотрим вариационный подход \cite{Bishop2006} для решения этой задачи. Пусть задано распределение $q$:
\begin{equation} \label{task.6}
q(\mathbf{w}) \sim \Lap(\mbox{\boldmath$\mu$}, \mbox{\boldmath$\Sigma$}), 
\end{equation}
где $\mbox{\boldmath$\mu$}, \mbox{\boldmath$\Sigma$}$ -- вектор средних и матрица ковариации, апроксимирующие неизвестное апостериорное распределение $p(\mathbf{w} \mid \mathfrak{D}, \mathcal{A})$, полученное при априорном предположении:
\begin{equation} \label{task.7}
p(\mathbf{w} \mid \mathcal{A}) \sim \Lap(\mathbf{m}, \mbox{\boldmath$\Theta$}),
\end{equation}
где $\mathbf{m}, \mbox{\boldmath$\Theta$}$ -- вектор средних и матрица ковариации априорного распределения.

Приблизим интеграл \ref{task.5} методом, предложенным в \cite{Bishop2006}:
\[
\label{task.8}
\begin{aligned}
\mathcal{L}_{\mathcal{A}}(\mathfrak{D}, \mathcal{A}) = \log p(\mathfrak{D} \mid \mathcal{A}) = 
\end{aligned}
\]

\[
\begin{aligned}
= \int_{\mathbf{w \in \mathbb{W}_{\mathcal{J}}}} q(\mathbf{w}) \log \frac{p(\mathfrak{D}, \mathbf{w} \mid \mathcal{A})}{q(\mathbf{w})} d \mathbf{w} - \int_{\mathbf{w \in \mathbb{W}_{\mathcal{J}}}} q(\mathbf{w}) \log \frac{p(\mathbf{w} \mid \mathfrak{D}, \mathcal{A})}{q(\mathbf{w})} d \mathbf{w} \approx 
\end{aligned}
\]

\[
\begin{aligned}
\approx \int_{\mathbf{w \in \mathbb{W}_{\mathcal{J}}}} q(\mathbf{w}) \log \frac{p(\mathfrak{D}, \mathbf{w} \mid \mathcal{A})}{q(\mathbf{w})} d \mathbf{w} =
\end{aligned}
\]

\[
\begin{aligned}
= \int_{\mathbf{w \in \mathbb{W}_{\mathcal{J}}}} q(\mathbf{w}) \log \frac{p(\mathbf{w} \mid \mathcal{A})}{q(\mathbf{w})} d \mathbf{w} + \int_{\mathbf{w \in \mathbb{W}_{\mathcal{J}}}} q(\mathbf{w}) \log p(\mathfrak{D} \mid \mathbf{w}, \mathcal{A}) d \mathbf{w} = 
\end{aligned}
\]

\[
\begin{aligned}
= \mathcal{L}_{\mathbf{w}}(\mathfrak{D}, \mathcal{A}, \mathbf{w}) + \mathcal{L}_{E}(\mathfrak{D}, \mathcal{A}).
\end{aligned}
\]

Первое слагаемое формулы \ref{task.8} -- это сложность модели. Оно определяется расстоянием Кульбака-Лейблера:
\begin{equation} \label{task.9}
\mathcal{L}_{\mathbf{w}}(\mathfrak{D}, \mathcal{A}, \mathbf{w}) = -D_{\KL}(q(\mathbf{w}) \parallel p(\mathbf{w} \mid \mathcal{A})).
\end{equation}
Второе слагаемое формулы \ref{task.8} является матожиданием правдоподобия выборки $\mathcal{L}_{\mathfrak{D}}(\mathfrak{D}, \mathcal{A}, \mathbf{w})$. В данной работе оно является функцией ошибки:
\begin{equation} \label{task.10}
\mathcal{L}_{E}(\mathfrak{D}, \mathcal{A}) = \mathsf{E}_{\mathbf{w} \sim q} \mathcal{L}_{\mathfrak{D}}(\mathbf{y}, \mathfrak{D}, \mathcal{A}, \mathbf{w}).
\end{equation}

Требуется найти параметры, доставляющие минимум суммарному функционалу потерь $\mathcal{L}_{\mathcal{A}}(\mathfrak{D}, \mathcal{A})$ из \ref{task.8}:
\begin{equation} \label{task.11}
\hat{\mathbf{w}} = \underset{\mathcal{A} \subset \mathcal{J}, \mathbf{w \in \mathbb{W}_{\mathcal{A}}}} \argmin - \mathcal{L}_{\mathcal{A}}(\mathfrak{D}, \mathcal{A}) = \underset{\mathcal{A} \subset \mathcal{J}, \mathbf{w \in \mathbb{W}_{\mathcal{A}}}} \argmin D_{\KL}(q(\mathbf{w}) \parallel p(\mathbf{w} \mid \mathcal{A})) - \mathcal{L}_{\mathfrak{D}}(\mathfrak{D}, \mathcal{A}, \mathbf{w}).
\end{equation}


% Методы решения задачи
\newpage

\section{Предлагаемый метод анализа распределения параметров нейросети}
В качестве априорного распределения параметров $\mathbf{w}$ рассмотрим $k$-мерное распределение Лапласа $\Lap$ с плотностью распределения:
\begin{equation} \label{methods.1}
f(\mathbf{x}) = \frac{2e^{\mathbf{x}^{\mathsf{T}}\mbox{\boldmath$\Sigma^{-1}$}\mbox{\boldmath$\mu$}}}{(2\pi)^{\frac{k}{2}}|\mbox{\boldmath$\Sigma$}|^{0,5}} \left( \frac{\mathbf{x}^{\mathsf{T}}\mbox{\boldmath$\Sigma^{-1}$}\mathbf{x}}{2 + \mbox{\boldmath$\mu$}^{\mathsf{T}}\mbox{\boldmath$\Sigma^{-1}$}\mbox{\boldmath$\mu$}} \right)^{\frac{v}{2}} K_v \left( \sqrt{(2 + \mbox{\boldmath$\mu$}^{\mathsf{T}}\mbox{\boldmath$\Sigma^{-1}$}\mbox{\boldmath$\mu$})(\mathbf{x}^{\mathsf{T}}\mbox{\boldmath$\Sigma^{-1}$}\mathbf{x})} \right),
\end{equation}
где $\mbox{\boldmath$\mu$} \in \mathbb{R}^k$ -- математичекое ожидание, $\mbox{\boldmath$\Sigma$} \in \mathbb{R}^{k \times k}$ -- матрица ковариации, $v = (2-k)/2$ и $K_v$ -- модифицированная функция Бесселя второго рода.

Рассмотрим случай, когда $k = 2$, $\mathbf{m} = \mathbf{0}$, а матрицы $\mbox{\boldmath$\Theta$}$ и $\mbox{\boldmath$\Sigma$}$ диаганальны. В таком случае плотность распределения $p(\mathbf{w} \mid \mathcal{A}) \sim \Lap(\mathbf{m}, \mbox{\boldmath$\Theta$}) \sim \Lap(\mathbf{0}, \mbox{\boldmath$\Theta$})$ будет иметь вид:
\begin{equation} \label{methods.2}
f(\mathbf{x}) = \frac{2}{(2\pi)^{\frac{k}{2}}|\mbox{\boldmath$\Theta$}|^{0,5}} K_0 \left( \sqrt{2 \mathbf{x}^{\mathsf{T}}\mbox{\boldmath$\Theta^{-1}$}\mathbf{x}} \right).
\end{equation}

Рассмотрим несколько случаев:
\begin{enumerate}
\item Пусть вектор параметров мал по норме, такой, что:
\begin{equation} \label{methods.3}
||\mathbf{w}|| \lesssim \frac{2}{||\mbox{\boldmath$\Sigma$}||}.
\end{equation}
В таком случае можно использовать аппроксимацию функции Бесселя \cite{Temme1975}:
\begin{equation} \label{methods.4}
K_v(x) \simeq \frac{\Gamma(v)}{2} \left( \frac{2}{x} \right)^v \Rightarrow K_0(x) \sim \frac{\Gamma(0)}{2} \sim \const.
\end{equation}
Тогда можно оценить:
\[
\label{methods.5}
\begin{aligned}
D_{\KL}(q(\mathbf{w}) \parallel p(\mathbf{w} \mid \mathcal{A})) = \int_{\mathbf{w \in \mathbb{W}_{\mathcal{J}}}} q(\mathbf{w}) \log \frac{p(\mathbf{w} \mid \mathcal{A})}{q(\mathbf{w})} d \mathbf{w} = 
\end{aligned}
\]

\[
\begin{aligned}
= \int_{\mathbf{w \in \mathbb{W}_{\mathcal{J}}}} q(\mathbf{w}) \log p(\mathbf{w} \mid \mathcal{A}) q(\mathbf{w}) d \mathbf{w} - \int_{\mathbf{w \in \mathbb{W}_{\mathcal{J}}}} q(\mathbf{w}) \log q(\mathbf{w}) d \mathbf{w} \sim
\end{aligned}
\]

\[
\begin{aligned}
\sim \int_{\mathbf{w \in \mathbb{W}_{\mathcal{J}}}} \frac{e^{\mathbf{w}^{\mathsf{T}}\mbox{\boldmath$\Sigma^{-1}$}\mbox{\boldmath$\mu$}}}{|\mbox{\boldmath$\Sigma$}|^{0,5}} K_0 \left( \sqrt{(2 + \mbox{\boldmath$\mu$}^{\mathsf{T}}\mbox{\boldmath$\Sigma^{-1}$}\mbox{\boldmath$\mu$})(\mathbf{w}^{\mathsf{T}}\mbox{\boldmath$\Sigma^{-1}$}\mathbf{w})} \right) \log \frac{1}{|\mbox{\boldmath$\Sigma$}|^{0,5}} \left(  K_0 \left( \sqrt{\mathbf{w}^{\mathsf{T}}\mbox{\boldmath$\Theta^{-1}$}\mathbf{w}} \right) \right) d \mathbf{w} -
\end{aligned}
\]

\[
\begin{aligned}
1
\end{aligned}
\]

\item Пусть норма вектора параметров большая:
\begin{equation} \label{methods.3}
||\mathbf{w}|| \sim ?????.
\end{equation}
Тогда функцию Бесселя можно аппроксимировать следующим образом [ссылка на https://dlmf.nist.gov/10.40]:
\begin{equation} \label{methods.3}
K_v(x) \simeq \left( \frac{\pi}{2x}\right)^{\frac{1}{2}}e^{-x} \sum \limits_{n=0}^{\infty} \frac{a_n(v)}{x^n} \Rightarrow K_0(x) \sim \frac{e^{-x}}{\sqrt{x}} \sum \limits_{n=0}^{10} \frac{a_n(0)}{x^n}.
\end{equation}
В этом случае можно получить следующую оценку:
\[1\] 
\end{enumerate}




% Вычислительный эксперимент
%\newpage

\section{Вычислительный эксперимент}
To be done...

% Выводы
%\newpage

\section{Заключение}
%\begin{table}[h!t]
%\begin{center}
%\caption{Результаты работы алгоритма}
%\label{table_2}
%\begin{tabular}{|c|c|c|c|c|}
%\hline
%	Ряд,~$\textbf{x}$ &Длина ряда,~$N$& \# сегментов,~$K$&Длина сегмента,~$T$& Ошибка,~$S$\\
%	\hline
%	\multicolumn{1}{|l|}{Phys.~Motion~1}
%	& 900& 2& 40& 0.06\\
%	\hline
%	\multicolumn{1}{|l|}{Phys.~Motion~2}
%	& 900& 2& 40& 0.03\\
%	\hline
%	\multicolumn{1}{|l|}{Synthetic~1}
%	& 2000& 2& 20& 0.04\\
%	\hline
%	\multicolumn{1}{|l|}{Synthetic~2}
%	& 2000& 3& 20& 0.03\\
%\hline

%\end{tabular}
%\end{center}
%\end{table}


В работе рассматривалась задача прореживания моделей нейросетей. Рассматривался метод оптимального прореживания и метод, основаный на вариационном подходе. Был предложен алгоритм прореживания, основаный на методе Белсли для удаления зависимых параметров модели. В ходе эксперимента было показано, что нейросети прорежены методом Белсли являются более устойчивы к шуму на входных данных. Качества прогноза нейросетей после прореживания методом Белсли не хуже качества прогноза нейросетей, полученных другими методами.


% Библиографические ссылки
\newpage


\begin{thebibliography}{99}
%	\bibitem{kwapisz2010}
%	\textit{J. R. Kwapisz, G. M. Weiss, S. A. Moore} Activity Recognition using Cell Phone Accelerometers~// Proceedings of the Fourth International Workshop on Knowledge Discovery from Sensor Data, 2010. Vol. 12. P. 74--82.
	
%	\bibitem{wang2014}
%	\textit{W. Wang, H. Liu, L. Yu, F. Sun} Activity Recognition using Cell Phone Accelerometers~// Joint Conference on Neural Networks, 2014. P. 1185--1190.
	
%	\bibitem{Ignatov2015}
%	\textit{A. D. Ignatov, V. V. Strijov} Human activity recognition using quasiperiodic time series collected from a single tri-axial accelerometer.~// Multimedial Tools and Applications, 2015.
	
%	\bibitem{Olivares2012}
%	\textit{A. Olivares, J. Ramirez, J. M. Gorris, G. Olivares, M. Damas} Detection of (in)activity periods in human body motion using inertial sensors: A comparative study.~// Sensors, 12(5):5791–5814, 2012.
	
%	\bibitem{cinar2018}
%	\textit{Y. G. Cinar and H. Mirisaee} Period-aware content attention RNNs for time series forecasting with missing values~// Neurocomputing, 2018. Vol. 312. P. 177--186.
	
%	\bibitem{motrenko2015}
%	\textit{A. P. Motrenko, V. V. Strijov} Extracting fundamental periods to segment biomedical signals~// Journal of Biomedical and Health Informatics, 2015,~20(6). P.~1466~-~1476.
	
%	\bibitem{lukashin2003}
%	\textit{Y. P. Lukashin} Adaptive methods for short-term forecasting~// Finansy and Statistik, 2003.
	
%	\bibitem{Ivkin2015}
%	\textit{И. П. Ивкин,  М. П. Кузнецов} Алгоритм классификации временных рядов акселерометра по комбинированному признаковому описанию.~// Машинное обучение и анализ данных, 2015.
	
%	\bibitem{Katrutsa2015}
%	\textit{V. V. Strijov, A. M. Katrutsa} Stresstes procedures for features selection algorithms.~// Schemometrics and Intelligent Laboratory System, 2015.
	
%	\bibitem{Borg2005}
%	\textit{I. Borg, P. J. F. Groenen} Modern Multidimensional Scaling. --- New York: Springer, 2005. 540 p.
	
%	\bibitem{Shiglavsi1997}
%	\textit{Д. Л. Данилова, А. А. Жигловский} Главные компоненты временных рядов: метод "Гусеница".~---~Санкт-Петербурскиий университет, 1997.
	
	\bibitem{Grabavoy}
	\textit{А.В. Грабовой, О.Ю. Бахтеев, В.В. Стрижов} Определение релевантности параметров нейросети.
	
	\bibitem{Bishop2006}
	\textit{Bishop C}. Pattern Regoghition and Machine Learninng. -- Berlin: Springer, 2006. 758p.
	
	\bibitem{Temme1975}
	\textit{Temme N}. On the numerical evaluation of the modified bessel function of the third kind~// Journal of Computational Physics, 1975. Vol. 19, 314p.
\end{thebibliography}

% Приложения
%\input{./appendices.tex}

\end{document}
